\section{Degree-four Salem parameters}\label{sal:section}

For a degree-four Salem number $\beta$, the two conjugates on the unit
circle form a single complex-conjugate pair.  We use simultaneous control
of this pair to construct many Fourier phases with appreciable expectation.
After removing integer trace terms, widely separated phases depend, up to a
small error, on disjoint blocks of the Bernoulli signs.  Their averages
therefore distinguish the Bernoulli convolution from Lebesgue measure.

\begin{theorem}\label{sal:main}
Let $\beta\in(1,2)$ be a Salem number of degree four.  Then the unbiased
Bernoulli convolution $\nu_{1/\beta}$ is singular with respect to Lebesgue
measure.
\end{theorem}

The proof has two quantitative requirements.  The Fourier products must
lose little mass before the algebraic trace becomes an accurate integer
approximation; the remaining trace errors must then stay small over many
disjoint blocks.  The multiplier construction below supplies both
requirements.  The trace approximation is the familiar arithmetic
ingredient in the Fourier analysis of Pisot and Salem parameters
\cite{Erdos1939}; see in particular Salem's small-conjugate construction
\cite[Chapter~III, Section~3, Theorem~III, pp.~28--29]{Salem1963}.
The estimates here also control the losses
accumulated while the multiplier changes.

\subsection{Multipliers with small unit-circle conjugates}

Fix $\beta$ as in Theorem~\ref{sal:main}, and write its conjugates as
$\beta^{-1},z,\bar z$, where $z=e^{i\theta}$.  The number $\beta$ is a unit:
the product of its four conjugates is~$1$.  Moreover, $\theta/\pi$ is
irrational, since otherwise $z$, and hence its conjugate $\beta$, would be
a root of unity.

For positive integers $k_1,k_2,\ldots$, to be chosen, define
\begin{equation}\label{sal:multipliers}
\begin{aligned}
 A_m&=\prod_{h=1}^m(\beta^{k_h}-\beta^{-k_h}),
 &b_m&=\sum_{h=1}^m k_h,\\
 c_m&=\prod_{h=1}^m(z^{k_h}-z^{-k_h}),
 &\delta_m&=|c_m|,\qquad L_m=\log(1/\delta_m).
\end{aligned}
\end{equation}
Here and throughout this section logarithms are natural.  Irrationality
ensures $\delta_m>0$, and $0<A_m\le\beta^{b_m}$.  For $j\in\Z$, put
$u_{m,j}=A_m\beta^j$.  All these numbers are algebraic integers, because
$\beta$ is a unit.  Their field traces are the integers
\begin{equation}\label{sal:trace}
 T_{m,j}=u_{m,j}+(-1)^m u_{m,-j}
                  +2\operatorname{Re}(c_mz^j)\in\Z.
\end{equation}
Indeed, the embedding taking $\beta$ to $\beta^{-1}$ takes $A_m$ to
$(-1)^mA_m$.  Also $u_{m,j}\notin\Z+\tfrac12$, since a rational algebraic
integer is an integer.  Thus none of the cosine factors used below
vanishes.

The construction below balances the decrease of $\delta_m$ against the
cosine losses generated by $A_m$.  Eventually $L_m$ grows by a factor
greater than two while $k_{m+1}\delta_m^2$ tends to zero; these are the
margins used in the loss estimate below.  It also keeps $b_m$ small
compared with $\delta_m^{-1}$, leaving room for many separated blocks of
signs.  A uniform lower bound on the cosines will permit a quadratic
estimate for perturbations of their logarithms.

\begin{lemma}\label{sal:construction}
There are positive integers $k_m$ and a constant $\gamma>0$ for which
the numbers $L_m$ in \eqref{sal:multipliers} are positive and satisfy
the following properties.
Set
\[
 D=\frac{10}{\log\beta},\qquad
 K_m=\left\lceil2b_m+DL_m\right\rceil\quad(m\ge2).
\]
Then $A_m\ge1$ for every $m$, and
\begin{align}
 K_m\le k_{m+1}\le e^{1.9L_m},\qquad
 L_{m+1}&\ge1.75L_m &&(m\ge2),\label{sal:growth-initial}\\
 \log K_m\le .75L_m,\qquad
 L_{m+1}&\ge2.1L_m &&(m\ge3).\label{sal:growth} 
\end{align}
Moreover,
\begin{equation}\label{sal:cosine-floor}
 |\cos(\pi u_{m,j})|\ge\gamma
       \quad(m\ge1,\ j\in\Z),
\end{equation}
and, for $m\ge2$,
\begin{equation}\label{sal:small-errors}
 2\delta_m\le\frac{\gamma}{2\pi},\qquad
 \beta^{b_m-k_{m+1}}\le e^{-10L_m}
                       \le\frac{\gamma}{2\pi}.
\end{equation}
\end{lemma}

\begin{proof}
Write $\|x\|_{\R/\Z}$ for distance to the nearest integer.  The pigeonhole
principle gives, for every integer $N\ge1$, an integer $q\in[1,N]$ such
that
\begin{equation}\label{sal:dirichlet}
 \|q\theta/\pi\|_{\R/\Z}\le N^{-1}.
\end{equation}
For example, place the $N+1$ fractional parts of
$0,\theta/\pi,\ldots,N\theta/\pi$ in $N$ intervals of length $1/N$.
Since $\theta/\pi$ is irrational, the possible $q$ in
\eqref{sal:dirichlet} are unbounded as $N\to\infty$.  In particular there
are arbitrarily large $q$ with $\|q\theta/\pi\|_{\R/\Z}\le1/q$.

Choose and fix $k_1$ so that $A_1>1$ and $\delta_1<1/16$.  Define
\[
 \gamma_m=\inf_{j\in\Z}|\cos(\pi u_{m,j})|.
\]
At this point $\gamma_1>0$: as $j\to-\infty$, $u_{1,j}\to0$; as
$j\to+\infty$, the trace identity shows that the absolute cosine differs
from $|\cos(\pi u_{1,-j})|$ by at most $2\pi\delta_1<\pi/8$.
The remaining finite set of factors contains no zero.

Next choose $k_2>k_1$ from the unbounded sequence satisfying
$\|k_2\theta/\pi\|_{\R/\Z}\le1/k_2$.  We will impose finitely many
lower thresholds on this choice.  Since
\[
 \delta_2\le\frac{2\pi\delta_1}{k_2},\qquad
 L_2\ge\log k_2-\log(2\pi\delta_1),
\]
both $k_2$ and $L_2$ can be made arbitrarily large.  In particular,
$\log K_2\le1.1L_2$ for every sufficiently large such choice: the last
inequality bounds $k_2$ by a fixed multiple of $e^{L_2}$, whereas the
other contribution $DL_2$ to $K_2$ is only linear in $L_2$.

For $m\ge2$, suppose $L=L_m$ is large and $\log K_m\le1.1L$.
Take $N=\lfloor e^{1.9L}\rfloor$ and choose $q$ by
\eqref{sal:dirichlet}.  Define
\begin{equation}\label{sal:choice-k}
 k_{m+1}=\begin{cases}
 q,&q\ge K_m,\\
 \lceil K_m/q\rceil q,&q<K_m.
 \end{cases}
\end{equation}
Let $L'=L_{m+1}$.  If $q\ge K_m$, then
\[
 \delta_{m+1}\le2\pi\delta_m/N,\qquad
 L'\ge2.9L-O(1)\ge2.8L,
 \qquad b_{m+1}\le2e^{1.9L}.
\]
If $q<K_m$, the multiplier $\lceil K_m/q\rceil$ is at most $2K_m$
and $K_m\le k_{m+1}<2K_m$.  Using
$|\sin(\pi x)|\le\pi\|x\|_{\R/\Z}$ gives
\[
 \delta_{m+1}\le4\pi\delta_mK_m/N,
 \qquad L'\ge2.9L-\log K_m-O(1)\ge1.75L,
\]
and $b_{m+1}\le3e^{1.1L}$.  In both cases
$K_m\le k_{m+1}\le e^{1.9L}$ when $L$ is sufficiently large.

These estimates also give $\log K_{m+1}\le .75L'$.  To see this
without an upper bound on $L'$, write
\[
 K_{m+1}\le2b_{m+1}+DL'+1.
\]
The bound on $b_{m+1}$ makes the first term at most
$\tfrac12e^{.75L'}$ for large $L$, since
$1.9<.75\cdot2.8$ in the first case and
$1.1<.75\cdot1.75$ in the second.  The remaining term satisfies the
same bound for every sufficiently large $L'$.  Thus the construction
iterates.  From $m=3$ onward, the improved bound
$\log K_m\le .75L_m$ gives
$L_{m+1}\ge(2.9-.75)L_m-O(1)\ge2.1L_m$ in the second case; the
first case already gives more.  This proves
\eqref{sal:growth-initial}--\eqref{sal:growth} after a lower threshold
on $L_2$ depending only on $\beta$.

It remains to choose $k_2$ large enough to obtain the cosine floor.
For $j\le0$ and $k=k_{m+1}$, the exact identity
\begin{equation}\label{sal:shift}
 u_{m+1,j}=u_{m,j+k}-u_{m,j-k},\qquad
 |u_{m,j-k}|\le\beta^{b_m-k}\beta^j
\end{equation}
and the Lipschitz bound for the cosine give a lower bound
$\gamma_m-\pi\beta^{b_m-k}$.
For positive $j$, reflect to $-j$ using \eqref{sal:trace} at level
$m+1$.  Consequently
\begin{equation}\label{sal:floor-recurrence}
 \gamma_{m+1}\ge\gamma_m
              -\pi\beta^{b_m-k_{m+1}}-2\pi\delta_{m+1}.
\end{equation}
For $m\ge2$, the inequality $k_{m+1}\ge2b_m+DL_m$ implies
$\beta^{b_m-k_{m+1}}\le e^{-10L_m}$.  Also
$L_m\ge1.75^{m-2}L_2$.  Thus the sum of the decrements in
\eqref{sal:floor-recurrence} is at most
\[
 \pi\beta^{k_1-k_2}+2\pi e^{-L_2}
 +\sum_{r=0}^{\infty}
   \left(\pi e^{-10\cdot1.75^rL_2}
         +2\pi e^{-1.75^{r+1}L_2}\right).
\]
This tends to zero as $k_2\to\infty$ along the chosen sequence.
Choose $k_2$ so that this sum is less than $\gamma_1/10$, and so that
$2e^{-L_2}$ and $e^{-10L_2}$ are at most $\gamma_1/(4\pi)$.
These requirements are compatible with the earlier lower thresholds.
Now fix $k_2$ and the resulting sequence, and put $\gamma=\gamma_1/2$.
Equations~\eqref{sal:cosine-floor}--\eqref{sal:small-errors} follow.
Finally, $k_m$ is increasing and every factor
$\beta^{k_m}-\beta^{-k_m}$ is at least its first factor $A_1>1$,
so $A_m\ge1$.
\end{proof}

\subsection{The loss in the cosine products}

Fix the sequence furnished by Lemma~\ref{sal:construction}.  Constants
denoted by $C$ below may change from line to line, but are independent of
$m$ and of the phase index.  Define
\[
 f(u)=-\log|\cos(\pi u)|
       \quad\text{for }u\notin\Z+\tfrac12,
 \qquad S_m=\sum_{j\le0}f(u_{m,j}).
\]
Each $S_m$ is finite: its terms have a geometrically decaying tail because
$u_{m,j}\to0$ as $j\to-\infty$ and $f(u)=O(u^2)$ near zero.
The next estimate controls the entire negative-power contribution even
though $A_m$ becomes large.

\begin{lemma}\label{sal:loss}
For the sequence in Lemma~\ref{sal:construction}, $S_m=o(L_m)$.
Furthermore, with $\eta=.01$, for $m\ge2$ and $n\ge0$,
\begin{equation}\label{sal:total-loss}
 \sum_{j\le n}f(u_{m,j})
       \le(2+\eta)S_m+Cn\delta_m^2.
\end{equation}
\end{lemma}

\begin{proof}
We first record a perturbation estimate.  If
$|\cos(\pi u)|\ge\gamma$ and $|v|\le\gamma/(2\pi)$, then
\begin{equation}\label{sal:taylor}
 f(u+v)\le(1+\eta)f(u)+Cv^2.
\end{equation}
Indeed, the absolute cosine is at least $\gamma/2$ along the segment
from $u$ to $u+v$, so $f''$ is bounded there.  Also
\[
 |f'(u)|\le\frac{\pi}{\gamma}\sqrt{2f(u)},
\]
because $1-|\cos(\pi u)|^2\le2f(u)$.  Taylor's formula and
$|f'(u)v|\le\eta f(u)+Cv^2$ prove \eqref{sal:taylor}.

The function $f$ is even and $1$-periodic.  Reflecting positive indices
by \eqref{sal:trace}, and using \eqref{sal:small-errors}, gives
\begin{equation}\label{sal:reflection-loss}
 f(u_{m,i})\le(1+\eta)f(u_{m,-i})+C\delta_m^2
                     \quad(m\ge2,\ i\ge1).
\end{equation}
Summing this inequality proves \eqref{sal:total-loss}.

To estimate $S_{m+1}$, apply \eqref{sal:taylor} to
\eqref{sal:shift}, with base point $u_{m,j+k}$ and
$v=-u_{m,j-k}$, where $k=k_{m+1}$.  The cosine floor holds at every
base point, including those with positive index.  The squared
perturbations have a geometric sum, so
\begin{align*}
 S_{m+1}
 &\le(1+\eta)\left(S_m+\sum_{i=1}^k f(u_{m,i})\right)
       +C\beta^{2(b_m-k)}\\
 &\le(1+\eta)(2+\eta)S_m+C(1+k\delta_m^2).
\end{align*}
By \eqref{sal:growth-initial},
$k\delta_m^2\le e^{-.1L_m}\le1$.
Therefore $S_{m+1}\le rS_m+C$, where
$r=(1.01)(2.01)=2.0301$.  The already fixed value $S_2$ is finite,
so iteration gives $S_m=O(r^{m-2})$.  In contrast,
\eqref{sal:growth} gives $L_m\ge2.1^{m-3}L_3$ for $m\ge3$.
Since $r<2.1$, their ratio tends to zero.
\end{proof}

\subsection{Independent blocks of Fourier phases}

We now turn the product estimates into sets of almost full
$\nu_{1/\beta}$-measure and vanishing Lebesgue measure.  The comparison
of phase averages under two measures is analogous to the separation
argument for Riesz products in \cite[Section~1.1, Theorem~1.2]{Peyriere1975}.
Here the needed estimates follow from the independent signs, as we now
show.  Work on the
probability space of independent uniform signs $(\epsilon_\ell)_{\ell\ge0}$,
and write
\[
 Y=\sum_{\ell\ge0}\epsilon_\ell\beta^{-\ell},\qquad
 W_{m,n}=\exp(i\pi A_m\beta^nY)\quad(n\ge0).
\]
The law of $Y$ is $\nu_{1/\beta}$, supported on
$I=[-(1-\beta^{-1})^{-1},(1-\beta^{-1})^{-1}]$.
Independence, followed by passage to the limit in the convergent sign
series, gives
\begin{equation}\label{sal:expectation}
 \E W_{m,n}=\prod_{j=-\infty}^n\cos(\pi u_{m,j}).
\end{equation}
This is a nonzero real number; the negative-index tail has a convergent
sum of logarithmic losses.

Set
\[
 J_m=\lfloor e^{1.5L_m}\rfloor,\qquad
 H_m=\lceil b_m+DL_m\rceil,
\]
\[
 \mathcal N_m=\{H_m+s(2H_m+1):s\in\Z_{\ge0}\}\cap[0,J_m],
\]
and let $M_m=|\mathcal N_m|$.  All subsequent estimates concern
sufficiently large $m$.  Since $H_m\le K_m\le e^{.75L_m}$,
\begin{equation}\label{sal:block-count}
 e^{.7L_m}\le M_m\le J_m+1.
\end{equation}
Lemma~\ref{sal:loss} and $J_m\delta_m^2\le e^{-.5L_m}$ show that
\begin{equation}\label{sal:mean-lower}
 |\E W_{m,n}|\ge e^{-L_m/20}
             \quad(0\le n\le J_m).
\end{equation}

For $n\in\mathcal N_m$, retain only the signs whose indices lie in
$[n-H_m,n+H_m]$, and define
\begin{equation}\label{sal:window-phase}
 \widetilde W_{m,n}
 =(-1)^{\sum_{j=H_m+1}^nT_{m,j}}
   \exp\left(i\pi\sum_{\ell=n-H_m}^{n+H_m}
                      u_{m,n-\ell}\epsilon_\ell\right).
\end{equation}
The integer sum in the prefactor is empty when $n=H_m$.
To explain this approximation, expand the exponent of $W_{m,n}$ as
$\pi\sum_{j\le n}u_{m,j}\epsilon_{n-j}$.  For $j>H_m$, replace
$u_{m,j}$ by its integer trace $T_{m,j}$; the resulting factor is
deterministic because
$e^{i\pi T_{m,j}\epsilon_{n-j}}=(-1)^{T_{m,j}}$.
For $j<-H_m$, drop the term.  The remaining indices are exactly those
in \eqref{sal:window-phase}.

The difference of the two phase angles, after these replacements, is
a sum of independent mean-zero signs.  The inequality
$|e^{ix}-e^{iy}|\le|x-y|$ therefore gives
\begin{align}
 \|W_{m,n}-\widetilde W_{m,n}\|_2^2
 &\le\pi^2\left(
       \sum_{j=H_m+1}^n|u_{m,j}-T_{m,j}|^2
                  +\sum_{j<-H_m}|u_{m,j}|^2\right)\notag\\
 &\le C\left(\beta^{2(b_m-H_m)}+n\delta_m^2\right)
 \le Ce^{-L_m/2}.\label{sal:window-error}
\end{align}
Here \eqref{sal:trace} bounds the first summand by a geometric tail
and a contribution $C n\delta_m^2$; the other sum is itself a geometric
tail.  Finally, $D\log\beta=10$ and $n\le J_m$ give the last inequality.
All norms in \eqref{sal:window-error} refer to the sign probability
space.  The windows belonging to consecutive members of $\mathcal N_m$
are disjoint, so the variables $\widetilde W_{m,n}$ are independent.

\begin{proof}[Proof of Theorem~\ref{sal:main}]
For $n\in\mathcal N_m$, let $\sigma_{m,n}\in\{-1,1\}$ be the sign of
the real expectation in \eqref{sal:expectation}.  Define the continuous
function
\[
 F_m(y)=\frac1{M_m}\sum_{n\in\mathcal N_m}
                       \sigma_{m,n}e^{i\pi A_m\beta^ny}.
\]
Its mean at $Y$ is real and, by \eqref{sal:mean-lower},
$a_m:=\E F_m(Y)\ge e^{-L_m/20}$.
The corresponding average of the independent variables
$\sigma_{m,n}\widetilde W_{m,n}$ has centered $L^2$ norm at most
$M_m^{-1/2}$, since each variable has modulus one.
By the triangle inequality and \eqref{sal:window-error},
\begin{equation}\label{sal:concentration}
 \|F_m(Y)-a_m\|_2
       \le M_m^{-1/2}+Ce^{-L_m/4}.
\end{equation}
Subtracting the means of the approximation errors introduces at most
a factor two, included in $C$.

On the other hand, the same average is small in Lebesgue $L^2(I)$.
Write its frequencies in increasing order as
$\xi_1<\cdots<\xi_{M_m}$.  Since $A_m\ge1$ and their exponents are
distinct nonnegative integers,
$\xi_{r+1}-\xi_r\ge\pi(\beta-1)$.
Thus $|\xi_r-\xi_s|\ge\pi(\beta-1)|r-s|$ and
\[
 \left|\int_I e^{i(\xi_r-\xi_s)y}\,dy\right|
       \le\frac2{|\xi_r-\xi_s|}\quad(r\ne s).
\]
Expanding the square and summing by the rank difference yields
\begin{equation}\label{sal:lebesgue-average}
 \int_I|F_m(y)|^2\,dy
       \le C\frac{1+\log M_m}{M_m}.
\end{equation}

Consider the compact sets
\[
 E_m=\{y\in I:\operatorname{Re}F_m(y)\ge\tfrac12e^{-L_m/20}\}.
\]
Chebyshev's inequality, \eqref{sal:concentration}, and
\eqref{sal:block-count} imply
\[
 1-\nu_{1/\beta}(E_m)
 \le4e^{L_m/10}\left(M_m^{-1/2}+Ce^{-L_m/4}\right)^2
 \longrightarrow0.
\]
Likewise, writing $|E_m|$ for Lebesgue measure,
\[
 |E_m|\le4e^{L_m/10}\int_I|F_m(y)|^2\,dy
       \le C e^{L_m/10}\frac{1+\log M_m}{M_m}
       \longrightarrow0,
\]
because $M_m\ge e^{.7L_m}$ and $\log M_m\le1.5L_m+O(1)$.
Choose a subsequence $(m_r)$ along which
$\sum_r(|E_{m_r}|+1-\nu_{1/\beta}(E_{m_r}))<\infty$.
Then $E=\limsup_r E_{m_r}$ is Borel and Lebesgue null, since
$|\bigcup_{r\ge R}E_{m_r}|\le\sum_{r\ge R}|E_{m_r}|\to0$.
Moreover, the summability of the complementary probabilities shows
that $\nu_{1/\beta}$-almost every point lies in all sufficiently late
$E_{m_r}$, and hence in $E$.  Therefore $\nu_{1/\beta}(E)=1$.
\end{proof}

\subsection{Two explicit parameters}\label{sal:examples-section}

\begin{corollary}\label{sal:examples}
Let $\beta_1$ and $\beta_2$ be the respective real roots greater than one of
\[
 P_1(x)=x^4-x^3-x^2-x+1,
 \qquad P_2(x)=x^4-2x^3+x^2-2x+1.
\]
Then $\beta_1,\beta_2\in(1,2)$ and both $\nu_{1/\beta_1}$ and
$\nu_{1/\beta_2}$ are singular.
\end{corollary}

\begin{proof}
For $y=x+x^{-1}$, the equations $P_i(x)/x^2=0$ become, respectively,
\[
 y^2-y-3=0,\qquad y^2-2y-1=0.
\]
Their larger roots are
$y_1=(1+\sqrt{13})/2$ and $y_2=1+\sqrt2$, both in $(2,5/2)$.
For each $i$, the equation $x+x^{-1}=y_i$ has one root
$\beta_i\in(1,2)$ and its reciprocal.  The other quadratic root
$y_i'$ lies in $(-2,2)$, so the remaining two roots of $P_i$ form a
nonreal conjugate pair on the unit circle.

It remains to check that these four roots are algebraic conjugates.
The quadratic field $\mathbb Q(y_i)$ is real, and the discriminant
$y_i^2-4$ has negative image $(y_i')^2-4$ under its nontrivial
embedding.  It cannot be a square in that field, since the image of a
square under a real embedding is nonnegative.  Hence
$x^2-y_ix+1$ is irreducible over $\mathbb Q(y_i)$, and
$[\mathbb Q(\beta_i):\mathbb Q]=4$.  Thus $\beta_1$ and $\beta_2$
are degree-four Salem numbers, and Theorem~\ref{sal:main} applies.
\end{proof}
